% =============================================================
% 3.4 – Thuộc chương THIẾT KẾ HỆ THỐNG
% =============================================================

\subsection{Thiết kế Giao diện (UI/UX Design)}

Phần thiết kế giao diện xác định bố cục tổng thể và luồng tương tác người dùng trước khi bước vào giai đoạn hiện thực. Nhóm áp dụng phong cách \textit{Mobile-First}, ưu tiên trải nghiệm trên điện thoại thông minh với màn hình nhỏ. Các nguyên tắc thiết kế cốt lõi bao gồm: phân cấp thông tin rõ ràng (thông số quan trọng nhất hiển thị lớn nhất), màu sắc nhất quán theo từng chức năng (xanh dương cho điều hướng, tím cho phòng/tầng, xanh lá cho chế độ tự động), và tối thiểu hóa số bước để thực hiện một tác vụ.

\subsubsection{Bản vẽ Wireframe các màn hình chính}

\paragraph{Màn hình 1 – Dashboard (Tổng quan hệ thống)}

\begin{figure}[H]
    \centering
    \includegraphics[width=0.38\linewidth]{image/ui_dashboard.png}
    \vspace{10pt}
    \caption{Giao diện màn hình Dashboard – Tổng quan cảm biến và thống kê thiết bị}
    \label{fig:wireframe_dashboard}
\end{figure}

\paragraph{Màn hình 2 – Quản lý không gian (Danh sách tầng/phòng)}

\begin{figure}[H]
    \centering
    \includegraphics[width=0.38\linewidth]{image/ui_spaces.png}
    \vspace{10pt}
    \caption{Giao diện màn hình Quản lý không gian – Danh sách tầng và phòng}
    \label{fig:wireframe_spaces}
\end{figure}

\paragraph{Màn hình 3 – Chi tiết phòng và điều khiển thiết bị}

\begin{figure}[H]
    \centering
    \includegraphics[width=0.38\linewidth]{image/ui_room_detail.png}
    \vspace{10pt}
    \caption{Giao diện màn hình Chi tiết phòng – Giám sát cảm biến và điều khiển thiết bị}
    \label{fig:wireframe_room}
\end{figure}

\paragraph{Màn hình 4 – Quản lý lịch trình}

\begin{figure}[H]
    \centering
    \includegraphics[width=0.38\linewidth]{image/ui_schedules.png}
    \vspace{10pt}
    \caption{Giao diện màn hình Lịch trình – Danh sách và trạng thái các lịch trình tự động}
    \label{fig:wireframe_schedule}
\end{figure}

\subsubsection{Luồng điều hướng ứng dụng (Navigation Flow)}

Ứng dụng được tổ chức theo mô hình \textbf{Tab Bar Navigation} kết hợp với \textbf{Stack Navigation}. Thanh điều hướng phía dưới cung cấp 4 tab chính, mỗi tab dẫn vào một nhóm chức năng riêng biệt:

\begin{table}[H]
\centering
\caption{Cấu trúc điều hướng của ứng dụng Mobile}
\begin{tabular}{|l|l|p{7cm}|}
\hline
\textbf{Tab} & \textbf{Màn hình gốc} & \textbf{Màn hình con (Stack)} \\ \hline
Dashboard & DashboardPage & (không có màn hình con) \\ \hline
Nhà & Home (Danh sách tầng) & FloorDetail $\rightarrow$ RoomDetail $\rightarrow$ AddDevice \\ \hline
Lịch trình & ScheduleListPage & ScheduleFormPage, ScheduleEditPage \\ \hline
Khác & MorePage & ActionLogPage, AutoModeConfigPage, UserManagementPage \\ \hline
\end{tabular}
\end{table}

% =============================================================
% 4.3 – Thuộc chương HIỆN THỰC HỆ THỐNG
% =============================================================

\subsection{Phát triển Ứng dụng Mobile}

\subsubsection{Frontend: Công nghệ và cấu trúc}

\paragraph{Công nghệ sử dụng}

Frontend được xây dựng như một \textbf{Single Page Application (SPA)} với bộ công nghệ hiện đại:

\begin{table}[H]
\centering
\caption{Danh sách công nghệ sử dụng cho Frontend}
\begin{tabular}{|l|l|p{6.5cm}|}
\hline
\textbf{Thư viện / Framework} & \textbf{Phiên bản} & \textbf{Mục đích sử dụng} \\ \hline
React & 18.3 & Thư viện UI dựa trên component, quản lý Virtual DOM. \\ \hline
TypeScript & --- & Kiểm tra kiểu tĩnh, nâng cao độ tin cậy của code. \\ \hline
Vite & 6.3 & Build tool tốc độ cao, hỗ trợ Hot Module Replacement. \\ \hline
TailwindCSS & 4.1 & CSS utility-first, định kiểu trực tiếp trên JSX. \\ \hline
Radix UI & 1.x & Bộ headless component (Dialog, Switch, Slider, ...) chuẩn WAI-ARIA. \\ \hline
React Router & 7.13 & Định tuyến phía client với hỗ trợ nested routes. \\ \hline
Recharts & 2.15 & Thư viện biểu đồ SVG (LineChart, PieChart) dựa trên D3. \\ \hline
Axios & 1.6 & HTTP client để gọi REST API từ Backend. \\ \hline
React Hook Form & 7.55 & Quản lý trạng thái biểu mẫu với hiệu năng cao. \\ \hline
Lucide React & 0.487 & Bộ icon SVG nhất quán về phong cách. \\ \hline
\end{tabular}
\end{table}

\paragraph{Cấu trúc thư mục Frontend}

\begin{verbatim}
src/
├── app/
│   ├── components/         # Reusable UI components
│   │   ├── ui/             # shadcn/ui primitives (Button, Card, Badge, ...)
│   │   ├── DeviceCard.tsx  # Card điều khiển từng thiết bị
│   │   └── layout/
│   │       └── AppShell.tsx  # Layout chứa Tab Bar và Protected Route
│   ├── context/
│   │   ├── AuthContext.tsx     # Quản lý JWT và thông tin người dùng
│   │   ├── SpaceContext.tsx    # State tầng, phòng (Floor, Room)
│   │   ├── DeviceContext.tsx   # State thiết bị và chế độ Manual/Auto
│   │   └── ScheduleContext.tsx # State danh sách lịch trình
│   ├── pages/
│   │   ├── dashboard/     # DashboardPage – Biểu đồ và tổng quan
│   │   ├── Home.tsx       # Danh sách tầng
│   │   ├── FloorDetail.tsx   # Danh sách phòng của một tầng
│   │   ├── RoomDetail.tsx    # Chi tiết phòng, cảm biến, điều khiển
│   │   ├── AddDevice.tsx     # Thêm thiết bị vào phòng
│   │   ├── schedules/     # Danh sách, tạo, chỉnh sửa lịch trình
│   │   ├── logs/          # Nhật ký hành động hệ thống
│   │   ├── auto-mode/     # Cấu hình ngưỡng chế độ tự động
│   │   └── more/          # Trang "Khác" (xuất dữ liệu, admin, ...)
│   └── routes.ts          # Khai báo toàn bộ routes của ứng dụng
└── services/
    └── api.ts             # Các hàm gọi API tập trung
\end{verbatim}

\paragraph{Các màn hình chính và chức năng hiện thực}

\begin{enumerate}
    \item \textbf{DashboardPage:} Là màn hình trung tâm, bao gồm ba vùng thông tin: (1) Vùng \textit{Glance} hiển thị tức thời nhiệt độ, độ ẩm, ánh sáng cùng chiều xu hướng so với 1 giờ trước; (2) Biểu đồ \textit{Hôm nay vs Hôm qua} dạng LineChart cho phép người dùng quan sát biến động trong ngày; (3) Vùng \textit{Thời lượng hoạt động} dùng DonutChart thể hiện phần trăm thời gian quạt và đèn LED được bật trong ngày, so sánh với hôm qua.

    \item \textbf{RoomDetail (Chi tiết phòng):} Phân tách danh sách thiết bị thành hai nhóm – \textit{Cảm biến} (sensors) và \textit{Thiết bị điều khiển} (controls). Tại đây, \texttt{DeviceCard} cho Quạt hiển thị thanh trượt PWM (0–100\%), còn DeviceCard cho Đèn RGB hiển thị các nút lựa chọn mức độ (Tắt / Cam / Trắng). Khi hệ thống ở chế độ tự động, một thanh cảnh báo xanh lá xuất hiện để người dùng biết thiết bị đang tự điều chỉnh.

    \item \textbf{ScheduleListPage / ScheduleFormPage:} Cho phép người dùng tạo và quản lý lịch trình với đầy đủ thông tin: tên lịch trình, giờ thực thi (\texttt{HH:MM}), ngày lặp lại trong tuần (T2–CN), thiết bị áp dụng và tham số điều khiển. Người dùng có thể bật/tắt từng lịch qua nút toggle, hoặc xóa và chỉnh sửa thông qua icon tương ứng.
\end{enumerate}

\subsubsection{Backend và API: Kiến trúc và giao tiếp}

\paragraph{Công nghệ Backend}

\begin{table}[H]
\centering
\caption{Stack công nghệ Backend}
\begin{tabular}{|l|l|p{6.5cm}|}
\hline
\textbf{Thư viện} & \textbf{Phiên bản} & \textbf{Vai trò} \\ \hline
Node.js + Express & 5.2 & HTTP server, định tuyến RESTful API. \\ \hline
mssql & 12.5 & Kết nối và truy vấn SQL Server (RDBMS). \\ \hline
mqtt & 5.15 & MQTT Client để subscribe feed thời gian thực từ Adafruit IO. \\ \hline
axios & 1.13 & Gọi Adafruit IO REST API để publish và lấy giá trị feed. \\ \hline
bcrypt & 6.0 & Hash và xác thực mật khẩu người dùng. \\ \hline
exceljs & 4.4 & Xuất dữ liệu nhật ký ra file Excel. \\ \hline
dotenv & 17.3 & Quản lý biến môi trường (DB, AIO Key, Port). \\ \hline
\end{tabular}
\end{table}

\paragraph{Kiến trúc các dịch vụ Backend}

Backend được tổ chức theo mô hình phân tầng với 5 dịch vụ nền chạy song song cùng HTTP server:

\begin{enumerate}
    \item \textbf{AdafruitService:} Đóng gói toàn bộ giao tiếp với Adafruit IO REST API. Cung cấp hai phương thức cốt lõi là \texttt{sendData(feedName, value)} để publish lệnh điều khiển và \texttt{getLastValue(feedName)} để đọc trạng thái hiện tại của feed khi server khởi động lại.

    \item \textbf{AdafruitMqttService:} Thiết lập kết nối MQTT liên tục đến Broker Adafruit IO. Khi nhận được dữ liệu mới từ các feed cảm biến (\texttt{DA\_TEMP}, \texttt{DA\_HUMI}, \texttt{DA\_BRIG}) hoặc feed chế độ (\texttt{DA\_MODE}), dịch vụ này cập nhật trạng thái vào bộ nhớ trong của \texttt{DeviceController} ngay lập tức.

    \item \textbf{AdafruitSyncService:} Chạy định kỳ mỗi 30 giây để đọc giá trị cảm biến từ \texttt{DeviceController} và lưu vào SQL Server. Điều này đảm bảo dữ liệu lịch sử được bảo tồn phục vụ cho các biểu đồ trên Dashboard.

    \item \textbf{AutoModeService:} Khi hệ thống ở chế độ tự động (\texttt{DA\_MODE = 1}), dịch vụ này chạy vòng lặp mỗi 10 giây. Logic tính toán dựa trên các ngưỡng có thể cấu hình từ ứng dụng: tốc độ quạt được quyết định bởi nhiệt độ (3 ngưỡng: thấp/trung/cao) và độ ẩm (cộng thêm lực đẩy nếu vượt ngưỡng); mức đèn được quyết định bởi cường độ sáng môi trường.

    \item \textbf{ScheduleExecutorService:} Chạy vòng lặp mỗi 60 giây, truy vấn toàn bộ lịch trình đang hoạt động trong Database và kiểm tra xem thời gian hiện tại có khớp không. Cơ chế \textit{deduplication} bằng Map nội bộ đảm bảo mỗi lịch trình chỉ thực thi một lần trong cùng một phút.
\end{enumerate}

\paragraph{Danh sách API Endpoints}

\begin{table}[H]
\centering
\caption{Danh mục các RESTful API Endpoint của hệ thống}
\begin{tabular}{|l|l|p{6cm}|}
\hline
\textbf{Method} & \textbf{Endpoint} & \textbf{Chức năng} \\ \hline
POST & \texttt{/api/auth/login} & Xác thực tài khoản, trả về token phiên làm việc. \\ \hline
GET & \texttt{/api/spaces/floors/:homeId} & Lấy danh sách tầng của một ngôi nhà. \\ \hline
POST & \texttt{/api/spaces/floors} & Tạo tầng mới. \\ \hline
DELETE & \texttt{/api/spaces/floors/:id} & Xóa tầng và toàn bộ phòng con. \\ \hline
GET & \texttt{/api/spaces/rooms/floor/:floorId} & Lấy danh sách phòng theo tầng. \\ \hline
POST & \texttt{/api/spaces/rooms} & Tạo phòng mới trong một tầng. \\ \hline
GET & \texttt{/api/devices/state} & Lấy trạng thái hiện tại của tất cả thiết bị. \\ \hline
POST & \texttt{/api/devices/fan} & Gửi lệnh điều khiển tốc độ quạt (0–100). \\ \hline
POST & \texttt{/api/devices/led} & Gửi lệnh điều khiển mức đèn LED (0/1/2). \\ \hline
POST & \texttt{/api/devices/lcd} & Gửi chuỗi văn bản hiển thị lên LCD. \\ \hline
GET & \texttt{/api/schedules} & Lấy danh sách tất cả lịch trình. \\ \hline
POST & \texttt{/api/schedules} & Tạo lịch trình mới. \\ \hline
PUT & \texttt{/api/schedules/:id/toggle} & Bật/tắt trạng thái lịch trình. \\ \hline
DELETE & \texttt{/api/schedules/:id} & Xóa lịch trình. \\ \hline
GET & \texttt{/api/logs/activity} & Lấy nhật ký hoạt động phân trang. \\ \hline
POST & \texttt{/api/auto-mode/config} & Lưu cấu hình ngưỡng chế độ tự động. \\ \hline
GET & \texttt{/api/dashboard/sensor-trend} & Lấy xu hướng cảm biến (so sánh 1h). \\ \hline
GET & \texttt{/api/dashboard/daily-runtime} & Lấy thời lượng hoạt động hôm nay/hôm qua. \\ \hline
GET & \texttt{/api/export/logs} & Xuất nhật ký dưới dạng file Excel (.xlsx). \\ \hline
GET & \texttt{/api/health} & Kiểm tra trạng thái kết nối DB và Adafruit. \\ \hline
\end{tabular}
\end{table}

\subsection{Tích hợp hệ thống}

Phần này mô tả quy trình kết nối toàn diện giữa ba thành phần: \textbf{Ứng dụng Mobile (App)}, \textbf{Backend Server} và \textbf{Phần cứng Yolo:bit} thông qua \textbf{Adafruit IO}.

\subsubsection{Luồng điều khiển thiết bị (Manual Control Flow)}

Khi người dùng thao tác trực tiếp trên ứng dụng, luồng xử lý diễn ra theo 6 bước liên tiếp:

\begin{enumerate}
    \item Người dùng kéo thanh trượt hoặc nhấn nút trên giao diện \textbf{RoomDetail}.
    \item \textbf{Frontend} gọi hàm trong \texttt{DeviceContext}, hàm này thực hiện HTTP POST đến Backend (ví dụ: \texttt{/api/devices/fan} với body \texttt{\{"speed": 60\}}).
    \item \textbf{Backend} tiếp nhận yêu cầu tại \texttt{DeviceController}, kiểm tra khóa ưu tiên (priority lock), sau đó gọi \texttt{AdafruitService.sendData("DA\_FAN", 60)}.
    \item \textbf{Adafruit IO} nhận dữ liệu qua REST API, lưu giá trị vào feed \texttt{DA\_FAN} và phân phối đến tất cả subscriber.
    \item \textbf{Yolo:bit} đang subscribe feed \texttt{DA\_FAN} nhận được sự kiện MQTT, khối lệnh xử lý sự kiện lập tức điều chỉnh tín hiệu PWM tại chân P1.
    \item Đồng thời, \textbf{Backend} ghi nhận hành động vào \texttt{System Log} với nguồn kích hoạt là \texttt{USER}.
\end{enumerate}

\subsubsection{Luồng đọc dữ liệu cảm biến (Sensor Data Flow)}

Dữ liệu cảm biến được cập nhật liên tục theo chu kỳ:

\begin{enumerate}
    \item Mỗi 10 giây, \textbf{Yolo:bit} đọc giá trị từ DHT20 (nhiệt độ, độ ẩm) và cảm biến ánh sáng, publish lên các feed \texttt{DA\_TEMP}, \texttt{DA\_HUMI}, \texttt{DA\_BRIG}.
    \item \textbf{AdafruitMqttService} (Backend) đang subscribe các feed này, nhận dữ liệu theo thời gian thực và cập nhật vào bộ nhớ trong của \texttt{DeviceController}.
    \item \textbf{AdafruitSyncService} chạy mỗi 30 giây, lấy dữ liệu từ \texttt{DeviceController} và lưu vào bảng \texttt{SensorLogs} trên SQL Server.
    \item \textbf{Frontend} gọi API \texttt{/api/devices/state} định kỳ để lấy giá trị mới nhất hiển thị lên giao diện và gọi \texttt{/api/dashboard/sensor-dual-chart} để vẽ biểu đồ.
\end{enumerate}

\subsubsection{Luồng thực thi lịch trình (Schedule Execution Flow)}

\begin{enumerate}
    \item Người dùng tạo lịch trình qua \textbf{ScheduleFormPage}, dữ liệu được lưu vào bảng \texttt{Schedules} và \texttt{ScheduleActionDefs} trên SQL Server.
    \item Mỗi 60 giây, \textbf{ScheduleExecutorService} truy vấn toàn bộ lịch trình có trạng thái \textit{active} từ Database.
    \item Với mỗi lịch trình, dịch vụ kiểm tra: thời gian hiện tại có khớp với \texttt{Time} và ngày hiện tại có nằm trong \texttt{RepeatDay} không.
    \item Nếu khớp và chưa thực thi trong phút này, dịch vụ gửi lệnh điều khiển đến Adafruit, ghi log với nguồn \texttt{SCHEDULE}, và nếu có cấu hình \texttt{DurationMinutes}, đặt hẹn giờ tự tắt sau khoảng thời gian đó.
\end{enumerate}

\subsubsection{Luồng chế độ tự động (Auto Mode Flow)}

\begin{enumerate}
    \item Người dùng bật chế độ Auto trên Dashboard. Frontend POST đến \texttt{/api/devices/mode} với body \texttt{\{"mode": 1\}}.
    \item Backend gửi giá trị \texttt{1} lên feed \texttt{DA\_MODE}, Yolo:bit nhận và chuyển sang logic tự động tại chỗ.
    \item Đồng thời, \textbf{AutoModeService} phía Backend được kích hoạt. Mỗi 10 giây, dịch vụ tính toán tốc độ quạt và mức đèn dựa trên dữ liệu cảm biến trong bộ nhớ, sau đó publish lệnh lên Adafruit nếu giá trị thay đổi so với chu kỳ trước.
    \item Cơ chế \textbf{Priority Lock} đảm bảo: lệnh thủ công của người dùng (nguồn \texttt{USER}) sẽ khóa thiết bị đó khỏi AutoMode trong 5 phút, tránh xung đột điều khiển.
\end{enumerate}

% =============================================================
% CHAPTER 5 – ĐÁNH GIÁ VÀ KẾT LUẬN
% =============================================================

\section{ĐÁNH GIÁ VÀ KẾT LUẬN}

\subsection{Kết quả đạt được}

\subsubsection{So sánh với mục tiêu ban đầu}

Nhìn lại 3 nhóm mục tiêu được đặt ra tại Chương 1, hệ thống đã hoàn thành đầy đủ các hạng mục theo kế hoạch:

\begin{table}[H]
\centering
\caption{Đánh giá mức độ hoàn thành so với mục tiêu ban đầu}
\begin{tabular}{|l|p{5.5cm}|c|p{3cm}|}
\hline
\textbf{Nhóm} & \textbf{Mục tiêu} & \textbf{Kết quả} & \textbf{Ghi chú} \\ \hline
Phần cứng & Thiết kế và lập trình Yolo:bit thu thập cảm biến, điều khiển quạt, đèn RGB, LCD. & \textbf{Hoàn thành} & Sai số cảm biến $\leq 1$°C. \\ \hline
Truyền dẫn & Hiện thực giao thức MQTT qua Adafruit IO với WiFi. & \textbf{Hoàn thành} & Tỷ lệ gửi nhận $\geq 99\%$ trong điều kiện mạng ổn định. \\ \hline
Backend & Xây dựng Server Node.js, API kết nối Adafruit và SQL Server. & \textbf{Hoàn thành} & 19 API endpoint, 5 dịch vụ nền. \\ \hline
Frontend & Phát triển ứng dụng Mobile quản lý đa tầng, lịch trình, cảnh báo. & \textbf{Hoàn thành} & 10 màn hình chức năng. \\ \hline
Database & Thiết kế CSDL lưu trữ cấu trúc nhà, lịch sử vận hành. & \textbf{Hoàn thành} & 8 bảng quan hệ, cascade delete. \\ \hline
\end{tabular}
\end{table}

\subsubsection{Đánh giá theo các yêu cầu chức năng}

\begin{table}[H]
\centering
\caption{Mức độ đáp ứng các yêu cầu chức năng (FR) đề ra}
\begin{tabular}{|l|p{6cm}|c|}
\hline
\textbf{Mã yêu cầu} & \textbf{Mô tả} & \textbf{Đã đáp ứng} \\ \hline
FR1 & Định nghĩa sơ đồ nhà (Nhà → Tầng → Phòng) & \checkmark \\ \hline
FR2 & Đăng ký thiết bị (Quạt, Đèn RGB, LCD, Cảm biến) vào phòng & \checkmark \\ \hline
FR3 & Cập nhật dữ liệu cảm biến thời gian thực & \checkmark \\ \hline
FR4 & Điều khiển Quạt (PWM 0–100\%) và Đèn RGB (3 mức) & \checkmark \\ \hline
FR5 & Lập lịch trình tự động theo giờ/phút và ngày trong tuần & \checkmark \\ \hline
FR6 & Tự động hóa theo ngưỡng cảm biến (chế độ Auto) & \checkmark \\ \hline
FR7 & Ghi nhận nhật ký hành động và hoạt động hệ thống & \checkmark \\ \hline
FR8 & Xuất nhật ký dưới dạng file Excel & \checkmark \\ \hline
\end{tabular}
\end{table}

\subsubsection{Đánh giá theo yêu cầu phi chức năng}

\begin{table}[H]
\centering
\caption{Đo lường thực tế so với chỉ tiêu phi chức năng}
\begin{tabular}{|l|l|l|l|}
\hline
\textbf{Mã chỉ số} & \textbf{Chỉ tiêu} & \textbf{Thực đo} & \textbf{Đạt} \\ \hline
NFR-P1 & Điều khiển $\leq$ 2 giây & $\approx$ 1.2--1.8 giây & \checkmark \\ \hline
NFR-P2 & Cập nhật cảm biến $\leq$ 5 giây/lần & 10 giây/chu kỳ phần cứng & \texttildelow{} \\ \hline
NFR-R1 & Hệ thống hoạt động $\geq$ 95\% thời gian & $\geq$ 97\% (trong thử nghiệm) & \checkmark \\ \hline
NFR-T1 & Cảnh báo gửi $\leq$ 3 giây & $\approx$ 2 giây qua AutoMode & \checkmark \\ \hline
NFR-A1 & Sai số lịch trình $\leq$ 5 giây & $\leq$ 1 phút (độ phân giải kiểm tra) & \texttildelow{} \\ \hline
\end{tabular}
\end{table}

\textit{Chú thích: \texttildelow{} nghĩa là đạt gần đúng nhưng chưa tối ưu hoàn toàn – chi tiết tại mục 5.2.}

\subsection{Hạn chế}

Mặc dù hệ thống hoàn thành đầy đủ các tính năng cốt lõi, quá trình kiểm thử thực tế cho thấy một số hạn chế kỹ thuật cần được ghi nhận:

\begin{itemize}
    \item \textbf{Phụ thuộc vào kết nối WiFi:} Toàn bộ luồng điều khiển đều yêu cầu kết nối Internet ổn định. Khi mạng bị ngắt, ứng dụng không có cơ chế fallback nội bộ (local control), khiến người dùng mất khả năng điều khiển từ xa.

    \item \textbf{Giới hạn băng thông Adafruit IO (Free Tier):} Tài khoản miễn phí của Adafruit IO giới hạn 30 data point mỗi phút. Với chu kỳ publish 10 giây (3 feed cảm biến × 6 lần/phút = 18 data point/phút), hệ thống hoạt động an toàn với một phòng; tuy nhiên khi mở rộng lên nhiều phòng song song, nguy cơ vượt ngưỡng là hiện hữu.

    \item \textbf{Chu kỳ lịch trình tối thiểu là 1 phút:} \texttt{ScheduleExecutorService} kiểm tra mỗi 60 giây nên sai số thực thi thực tế có thể lên đến 59 giây – cao hơn chỉ tiêu NFR-A1 đề ra là 5 giây. Để đạt độ chính xác dưới giây, cần thay đổi cơ chế scheduling sang \textit{event-driven} với bộ đếm thời gian chính xác hơn.

    \item \textbf{Quản lý một ngôi nhà duy nhất:} Kiến trúc hiện tại của Adafruit feed không hỗ trợ đa nhà một cách tự nhiên do các feed đang được cố định tên (\texttt{DA\_TEMP}, \texttt{DA\_FAN}, ...). Mở rộng cho nhiều ngôi nhà đồng thời đòi hỏi thiết kế lại cơ chế ánh xạ feed.

    \item \textbf{Chưa có push notification:} Hệ thống cảnh báo hiện tại hoạt động theo cơ chế polling từ Frontend. Khi người dùng thoát ứng dụng, họ sẽ không nhận được thông báo dù cảm biến vượt ngưỡng. Tính năng này đòi hỏi tích hợp dịch vụ như Firebase Cloud Messaging (FCM).

    \item \textbf{Giao diện chưa tối ưu trên màn hình lớn:} Thiết kế Mobile-First (max-width: 28rem) làm cho giao diện hiển thị không tối ưu trên tablet hoặc desktop, nơi phần lớn diện tích màn hình bị bỏ trống.
\end{itemize}

\subsection{Hướng phát triển}

Dựa trên nền tảng hiện có, nhóm xác định một số hướng phát triển tiếp theo có tính khả thi và giá trị thực tiễn cao:

\begin{enumerate}
    \item \textbf{Mở rộng hệ sinh thái cảm biến:} Tích hợp thêm cảm biến khí CO/CO$_2$ phục vụ phát hiện rò rỉ khí gas, cảm biến PIR (chuyển động) để tự động bật đèn khi có người, và cảm biến cửa/cửa sổ từ trường để giám sát an ninh. Mỗi loại cảm biến sẽ được ánh xạ vào một loại thiết bị mới trong bảng \texttt{DeviceTypes}.

    \item \textbf{Triển khai MQTT Broker nội bộ (Local Broker):} Thay thế Adafruit IO bằng một MQTT Broker tự triển khai (ví dụ: Mosquitto) chạy trên Raspberry Pi hoặc router nội bộ. Điều này loại bỏ phụ thuộc vào Internet, giảm độ trễ xuống còn \textless{}100ms và không còn bị giới hạn băng thông.

    \item \textbf{Ứng dụng Mobile Native (React Native):} Chuyển đổi từ Web App sang React Native để có thể đóng gói và phân phối qua App Store / Google Play, đồng thời tận dụng các tính năng nền tảng như push notification (FCM), background sync, và widget màn hình chính.

    \item \textbf{Tích hợp trí tuệ nhân tạo trong quản lý năng lượng:} Áp dụng các mô hình học máy đơn giản (ví dụ: hồi quy tuyến tính) để phân tích dữ liệu lịch sử \texttt{SensorLogs} và \texttt{DeviceLogs}, từ đó đưa ra gợi ý tối ưu hóa lịch trình hoặc cảnh báo sớm về bất thường trong tiêu thụ điện năng.

    \item \textbf{Tích hợp điều khiển bằng giọng nói:} Kết nối hệ thống với các nền tảng như Google Assistant hoặc Amazon Alexa thông qua các webhook tương thích IFTTT, cho phép người dùng điều khiển thiết bị bằng lệnh thoại mà không cần mở ứng dụng.

    \item \textbf{Hỗ trợ đa người dùng và phân quyền:} Mở rộng mô hình bảo mật để một ngôi nhà có thể được chia sẻ cho nhiều tài khoản với các mức quyền khác nhau (chủ nhà, thành viên gia đình, khách), phục vụ nhu cầu thực tế của các hộ gia đình đông thành viên hoặc môi trường doanh nghiệp nhỏ.
\end{enumerate}
